\chapter{System Overview and\newline Design Approach}
The E-ShoeS (ESS) project combines four main features inside a standard running shoe shape: footstep kinetic energy harvesting, GPS positioning, AI control with an emergency SOS system, and active lighting. The design uses fundamental engineering concepts to select lightweight and durable materials like carbon fiber and titanium to protect internal electronics. First-principles calculations were used to estimate energy production and system power consumption, followed by component testing to ensure all parts operate smoothly together.

\section{Energy Harvesting and Storage System}
\subsection{Footstep Energy Harvesting Components}
To gain energy from footsteps, piezoelectric transducers which transform mechanical stress or pressure into electrical energy or the other way around are set up in the sole. Its principle works based on the piezoelectric effect reacting to mechanical stress~\cite{ahmed2024smart}. A Power Management IC (PMIC) manages, adjusts, and supplies power to key parts inside a device, including processors, memory, and peripherals~\cite{jlcpcb2025pmic}.

\subsection{Energy Storage and Power Management Components}
A rechargeable Li-Po battery stores energy collected from footstep harvesters to power the microcontrollers, sensors, and lights. A power management IC regulates charging current and protects the battery from overcharging or deep discharging. A DC-DC voltage regulator adjusts battery output to a stable 3.3 V required by the BLE microcontroller and LED drivers.

\begin{table}[H]
    \centering
    \caption{Power subsystem component specifications.}
    \vspace{4pt}
    % \label{tab:energy_components}
    \begin{tabularx}{\textwidth}{L L L L L}
        \toprule
        \textbf{Component} & \textbf{Type} & \textbf{In/Output} & \textbf{Size} & \textbf{Position} \\
        \midrule
        Piezoelectric transducers & PolyK PVDF Film & Higher than 100~V AC output from steps & 20 $\times$ 67~mm (110~µm thick) & Embedded beneath the insole heel/forefoot \\
        \addlinespace
        PMIC & e-peas AEM30940 & Input: 50~mV--4.5~V;\newline Output: 1.8--4.1~V DC & 5 $\times$ 5~mm & Main electronic control PCB in midsole \\
        \addlinespace
        Rechargeable Battery & Li-Po 3.7~V 4000~mAh & 3.7~V nominal, 14.8~Wh capacity & 12 $\times$ 40 $\times$ 65~mm & Sealed compartment near shoe arch\\
        \addlinespace
        DC-DC Voltage Regulator & TPS63031 & Input: 1.8--5.5~V;\newline Output: 3.3~V DC & 2.5 $\times$ 2.5~mm & Main PCB, near microcontroller \\
        \bottomrule
    \end{tabularx}
\end{table}

\section{SOS system}
The EShoeS will trigger an SOS system when the user presses a button designed into the footwear. The button activating an emergency signal is designed to be an all in one system with the light-activation button, which is familiar with the power activation button on smartphones. The proposed design uses a press-and-hold mechanism, and users have to hold the button for approximately 2--3 seconds before triggering an emergency signal.

To send the signals from the shoes to the user's smartphone, Bluetooth Low Energy (BLE) is selected as the primary communication technology. BLE operates in the 2.4 GHz ISM band~\cite{bluetooth_sig_2023} and is known for its ultra-low power consumption~\cite{dementyev2013power}. Compared with direct cellular connectivity, it also offers a more compact design and is more affordable. These characteristics make it suitable for a wearable product with limited battery capacity. For this system, the MDBT42Q-512KV2 is chosen as the BLE module~\cite{raytac_mdbt42q}. It already combines the BLE microcontroller and a built-in chip antenna in one small unit to process and send the SOS signal to the smartphone. Inside this module, the nRF52832 microcontroller uses a 64 MHz Cortex-M4F core with a floating-point unit, offering enough computing performance and DSP capabilities~\cite{nordic_nrf52832}. 

The dedicated smartphone application acts as the receiver for emergency signals transmitted via BLE\@. Upon receiving an alert from the shoes, the app accesses the smartphone's built-in Global Navigation Satellite System (GNSS/GPS) to determine the user's exact geographical location and immediately dispatches emergency notifications containing these coordinates to pre-registered contacts over an active internet connection.

The internal rechargeable battery serves as the central power source for all on-shoe electronics, supplying energy to the SOS control module, BLE microcontroller, front headlamp, and RGB LED strip. Additionally, energy harvested from footsteps is directed to this main battery, ensuring continuous operation even when external power sources are unavailable.
